He had carried it for weeks now, since Dragonstone, since the night he’d sat across from a boy in a chair who was no longer entirely a boy and heard him say, in that same flat, faraway voice he used for everything now, \textit{his name is Aegon Targaryen, and he is the rightful heir to the Iron Throne} — and Sam had thought, at the time, that the knowing would be the hard part.

He understood now that he’d been wrong. The knowing had been easy. It was the carrying that was killing him.

He sat that evening in the small chamber Gilly had made warm and human out of stone that had never been built with warmth in mind, little Sam asleep at last after an hour of fussing that had left them both hollow-eyed, and found himself turning the same three facts over and over the way a man worries a stone in his pocket long after his hand has learned every edge of it. Rhaegar Targaryen had annulled his marriage to Elia Martell before he ever rode north. Rhaegar Targaryen had married Lyanna Stark in secret, beneath a godswood far from anyone who might have stopped it. And the boy born of that marriage, hidden his whole life beneath a bastard’s name because the truth of him would once have meant a swift and certain death, was currently somewhere above Sam’s head in this same castle, bracing for a war that might kill him before anyone ever told him who he actually was.

“You’re doing the thing again,” Gilly said, not looking up from the shirt she was mending, “where you go somewhere else behind your own eyes and forget I’m sitting across from you.”

“I’m sorry. I’ve a great deal to think through.”

“You’ve had a great deal to think through for weeks now, and you still haven’t told him.” She set the mending down, and for a moment she seemed to be hunting for the right shape of the thing rather than already holding it ready. “I don’t know how to say this so it comes out kind, Sam, so I’ll just say it plain and you can forgive me the roughness of it after. I know what it cost you, learning about your father. About Dickon.” She did not say \textit{burned}, though the word sat in the room with them regardless, the way it always did whenever that particular grief came near the surface. “I know that’s made you frightened of what truth costs, when it’s a queen with dragons deciding the price of it. But this isn’t — “ She stopped, frustrated with her own sentence, and started again. “Look. You’re not keeping it from him to spare him. You’re keeping it because saying it out loud makes it real, and once it’s real you can’t take it back and go on being just his friend the way you’ve always been. That’s not the same as protecting him, Sam, whatever you’ve been telling yourself it is. That’s protecting the two of you the way things are right now. He doesn’t get a vote in that, and I don’t think that’s fair to him, and I don’t much care whether it’s fair to you.”

Sam had turned some version of that argument over in his own mind a hundred times these past weeks, in every possible order, and heard it now, in Gilly’s rougher, less patient words, land somewhere it hadn’t quite managed to land before — not because she’d found the perfect sentence for it, he understood, but because she plainly hadn’t, and had said the true thing anyway rather than waiting for it to come out well. His father had died refusing to kneel to a claim he judged unworthy. It seemed a cruel, precise sort of irony that Sam alone, of everyone breathing in this castle, held the knowledge that might have changed which claim Randyll Tarly had been asked to kneel to in the first place — for all the good that irony did either of them now.

“I keep thinking,” he said slowly, “that if I tell him, I can’t ever untell him. That whatever peace he’s found, bending his knee, loving her the way I think he does — I’ll have broken something in him that doesn’t mend back the way it was.”

“Some things aren’t meant to mend back the way they were,” Gilly said. “Some things are only meant to be true, and left for the person they belong to, to carry however they choose to carry it.” She reached over, took his hand, callused and ink-stained both. “You’ve carried enough alone, Sam. Let him carry his own truth now. That’s not cruelty. That’s just giving a man back what’s his.”

\sceneBreak

He found Jon in the library, which did not surprise him overmuch. Sam had come to understand, these past weeks, that Jon had taken to sitting there most nights, long after the rest of the castle had banked its fires — ostensibly working through what little remained of the North’s own campaign histories, though Sam suspected, watching him some evenings from the doorway before he worked up the nerve to enter, that the reading mattered rather less to Jon than simply having somewhere quiet and unpeopled left to sit that wasn’t a war council or a stranger’s expectant face.

He was bent over a table cluttered with borrowed ledgers and half-legible campaign accounts when Sam found him tonight, a single candle burned low beside him, and he looked up with the particular weary attention of a man who has been reading for hours without absorbing much of anything.

“Trying to learn something useful about how the last Long Night was actually fought,” Jon said, before Sam had properly crossed the room, “since half the men in that hall seem to think shouting louder wins wars on its own. Not much luck so far. Whoever wrote most of these was more interested in flattering some dead lord’s ancestors than recording anything a soldier could use a thousand years later.”

“Jon.” Sam’s voice came out rougher than he’d meant it to, and he made himself go on before his own courage failed him the way it had failed him a dozen times already these past weeks. “There’s something I have to tell you. Something I should have told you a long while ago.”

Jon set the campaign history aside, and Sam saw, in his old friend’s face, the particular stillness of a man bracing for grief he already half-expected. “Is it Ned Stark? Something more about the rebellion—“

“You’re not his son.”

It came out of him wrong — too fast, too bare, no shape to it at all, and Sam heard himself say it and understood, a half-second too late, that he had just handed Jon the ending of a story without ever giving him the beginning of it, and Jon’s face did something confused and hurt at once, as though Sam had accused him of something rather than told him anything.

“I know I’m not his trueborn son, Sam. I’ve known that my whole—“

“No. That’s not — I’ve said it backward. I’m sorry.” Sam pressed the heels of his hands briefly against his own eyes, furious with himself, because he had rehearsed this a dozen times walking up the stairs and every rehearsed version had come out in the right order and none of that order had survived the actual saying of it. “Let me start again. Properly, this time.”

He made himself breathe.

“Bran and I found a record at Dragonstone, before you ever sailed north. An annulment — Rhaegar Targaryen’s marriage to Elia Martell, dissolved by the High Septon’s own hand. And then a second record, after it. A marriage. Secret, at a godswood, between Rhaegar and a woman the record names only as a Northern lady.” He stopped again, hearing himself, and forced the next part out before he could lose his nerve a second time. “Jon, your mother wasn’t some woman who caught a prince’s eye and died of it. Your mother was Lyanna Stark. I should have said that first. I don’t know why I didn’t say that first, it’s the whole of it, it’s the only part that actually matters to you and I buried it under an annulment like a fool.”

“Sam.” Jon’s voice had gone very quiet. “Say the rest of it properly, then. All of it, in the right order, and I’ll not interrupt you again till you’re done.”

Sam nodded, and made himself hold to that promise, though his hands had not stopped shaking.

“She died giving birth to you. At the Tower of Joy. Not of any nameless illness — Bran’s seen it, actually seen it, the way he sees things now, not pieced together from old records — and with her last breath she made your father promise her something, and the promise was to protect you from a mad king’s daughter and a usurper’s justice both, by raising you as his own bastard and never once telling you the truth. Not even when it might have been safe to.” His voice had gone thick, and he did not try to steady it. “You’re not Ned Stark’s bastard son, Jon. You never were. Your name is Aegon Targaryen, trueborn son of Rhaegar and Lyanna, and by every law of succession this country’s ever recognized, you stand before Daenerys Targaryen in the line to that throne, not behind her.”

The silence that followed was long enough that Sam heard, somewhere beyond the library’s shuttered windows, the muffled sound of the castle’s night watch changing over, entirely unaware that the ground had just shifted beneath one man sitting quietly among borrowed books.

“Say that again,” Jon said finally, very low. “Say it plainly, Sam. I need to hear it said plainly, not in a riddle.”

“You are the trueborn king of Westeros, Jon. Aegon Targaryen, sixth of that name to bear it, though I don’t imagine that particular numbering matters much to you right now. Daenerys is your aunt. And she has no idea, so far as either Bran or I have ever told her, that the man she’s fallen in love with holds a truer claim to her own throne than she does.”

Jon’s hand found the table’s edge, steadying himself against it the way a man steadies himself against a wall after a blow he didn’t see coming, and when he spoke again his voice had gone rough in a way Sam had rarely heard from him, not even in the worst of what they’d survived together at the Wall.

“Who else knows this.”

“Bran. Myself. No one else, so far as I’m aware — Bran said the record itself may not have survived, and there’s no one living who witnessed the marriage who’d dare speak of it now even if it had.”

“Then no one else will know it,” Jon said, and there was something in the flatness of it that unsettled Sam more than shouting would have. “Not Daenerys. Not Sansa, not Arya, not Tyrion, not one soul in this castle beyond the three of us who already carry it. I’ve just watched a hall full of Northern lords nearly come apart over whether I was right to bend my knee to her at all, Sam. I’ll not hand them a truth that says I never should have knelt in the first place, not while there’s a war to survive that needs both of us whole and united to win it. Whatever this makes me — whatever claim it gives me — it can wait. It has to wait. Some truths cost more told too soon than they’re worth.”

“And after the war?”

Jon did not answer that, and Sam, watching his oldest friend’s face in the library’s low candlelight, understood that he had just witnessed a man choose, deliberately and with his eyes open, to carry the heaviest kind of silence there was. Whether he had done Jon a kindness in giving him that choice, or handed him a burden crueler than the bastard’s name he’d carried his whole life, Sam had no way of knowing, and suspected the not-knowing would sit with him a long while yet.